
Predicting properties of trained neural networks from their weights --- model zoo analysis --- typically assumes that test models share the same architecture as training models. When this assumption breaks, flat tokenization methods fail: SANE, a state-of-the-art within-architecture weight-space SSL method, maps all ViT inputs to near-identical latent codes after training on CNN checkpoints, achieving $R^{2}=0.072$ (near-chance) on 53 held-out ViT-S/16 ImageNet checkpoints. We identify the root cause as the representation substrate rather than the learning algorithm: flat weight tokenization discards the relational graph structure that is consistent across architecture families. We train a permutation-equivariant graph SSL encoder (EquiSSL-perm) on MLP/CNN model zoos and evaluate zero-shot cross-architecture transfer to ViT property prediction, achieving $R^{2}=0.231$ versus SANE's $R^{2}=0.072$ --- a ~220\% relative improvement (3.$19\times$ ratio) --- without any ViT training data. Counterintuitively, adding scale equivariance to the symmetry group reduces performance: EquiSSL (scale+permutation) achieves $R^{2}=0.185$--0.210 depending on the evaluation setup, compared to EquiSSL-perm's $R^{2}=0.231$--0.327. We attribute this reversal to a LayerNorm gauge-fixing hypothesis: LayerNorm eliminates the scale degree of freedom that scale equivariance targets, making scale invariance an incorrect inductive bias for ViT weight encoding. All results are from a single training seed; statistical significance is not assessable. Latent interpolation fails entirely (acc=0.100, random chance), attributable to a decoder design limitation rather than latent space quality.

